% P08 | 3 Part A: Brownian motion and random walks | Result: sample paths
% guide: dev_notes/paragraph_guides/P08_result_sample_paths.txt
Figure~\ref{fig:part1} shows that the walk with $N_2=10$ makes very few turns and looks like a polygon, whereas the walk with $N_2=1000$ is hard to tell from a Brownian path.
Between these two the number of turns and the roughness of the path grow steadily with $N_2$.
The maxima and minima of the walks are comparable for every $N_2$, because the variance at $t=1$ is equal to $1$ whatever the step size.
What changes with $N_2$ is the fine structure of the path, not its overall range.
Neither object has a derivative, so the Brownian paths are not smoother than the walks: it is the walk that becomes Brownian as $N_2$ grows.
This is the point of Donsker's theorem.
A picture is suggestive but not a measurement, so the next three paragraphs quantify the convergence.
